% Source: https://artofproblemsolving.com/community/c3h1635553_sierpinski_triangle
% Stable identity: aops-topic-1635553-q01
Joining the midpoints of an equilateral triangle of side $s$ splits it into four equilateral triangles of side $\dfrac {s}{2}$ (each midsegment is half the parallel side). The first four steps look like this:

\begin{asy}
usepackage("amsmath");
size(9cm);
defaultpen(fontsize(10pt));
void sier(pair a,pair b,pair c,int n){
 if(n==0){fill(a--b--c--cycle,paleblue);draw(a--b--c--cycle,blue+linewidth(0.4));return;}
 sier(a,(a+b)/2,(a+c)/2,n-1);sier((a+b)/2,b,(b+c)/2,n-1);sier((a+c)/2,(b+c)/2,c,n-1);}
real h=sqrt(3)/2;
for(int k=0;k<4;++k){pair o=(1.25*k,0);sier(o,o+(1,0),o+(0.5,h),k);draw(o--o+(1,0)--o+(0.5,h)--cycle,linewidth(0.7));label("step "+string(k),o+(0.5,0),S,fontsize(9pt));}
label("\mbox{}",(min(currentpicture,user=true)+max(currentpicture,user=true))/2,black); // keeps the text after the figure black when LaTeX typesets the labels inline
\end{asy}

\begin{enumerate}
\item[(a)] Each step replaces every shaded triangle by $3$ shaded triangles of half the side. So after step $n$ there are $\boxed{3^n}$ shaded triangles, each of side $\boxed{\dfrac{1}{2^n}}$.
\item[(b)] Each step keeps $3$ of the $4$ equal parts of every shaded triangle, i.e.\ multiplies the shaded area by $\dfrac{3}{4}$. After step $n$ the shaded fraction is $\boxed{\left(\dfrac{3}{4}\right)^n}$. (Equivalently, $3^n$ triangles each with $\dfrac{1}{4^n}$ of the original area.)
\item[(c)] The total perimeter is $3^n\cdot3\cdot\dfrac{1}{2^n}=\boxed{3\left(\dfrac{3}{2}\right)^n}$.
\item[(d)] Since $0<\dfrac{3}{4}<1$, the powers $\left(\dfrac{3}{4}\right)^n$ tend to $0$: for example $\left(\dfrac{3}{4}\right)^{4}=\dfrac{81}{256}<\dfrac{1}{3}$, so every $4$ steps divide the shaded area by more than $3$, and after $4k$ steps the fraction is below $\dfrac{1}{3^k}$, which is eventually smaller than any given positive number. On the other hand $\left(\dfrac{3}{2}\right)^n\ge1+\dfrac {n}{2}$ (Bernoulli's inequality, or by induction: $\left(\dfrac{3}{2}\right)^{n+1}\ge\left(1+\dfrac {n}{2}\right)\dfrac{3}{2}\ge1+\dfrac{n+1}{2}$), so the total perimeter $3\left(\dfrac{3}{2}\right)^n\ge3+\dfrac{3n}{2}$ exceeds any given number for large $n$. So the shaded region gets arbitrarily small in area while its boundary gets arbitrarily long.
\end{enumerate}
